%% kapa: exam=1 round=21 year=2010 period=1 form=A subject=economics subject_ko=경제원론 seq=38 number=78
%% source: Q-Net 감정평가사 기출문제 (cst00309) artlSeq=5024181 file=extracted/1008/5024181/2010년 제21회 감정평가사 1차시험 기출문제/1교시.hwp
%% conversion: kordoc 4.24.51 -> tools/gen1.py (auto); review=proofread-sonnet; flags=none
\begin{question}{78}
생산요소시장과 생산물시장이 모두 완전경쟁적이다. 이에 관한 설명으로 옳지 않은 것은?
\begin{choices}
  \item 완전경쟁 기업의 이윤극대화를 위한 조건은 ‘생산요소의 한계생산물가치 = 한계요소비용’ 이다.
  \item 완전경쟁 기업의 이윤극대화를 위한 조건은 ‘생산요소의 한계수입생산물 = 생산요소의 가격’ 이다.
  \item 생산물시장이 완전경쟁적일 경우가 독점적일 경우보다 기업의 생산요소 수요가 더 작다.
  \item 생산요소의 한계생산곡선이 우측으로 이동하면 완전경쟁 기업의 생산요소 수요가 증가된다.
  \item 생산물에 대한 수요의 감소는 완전경쟁 기업의 생산요소 수요를 감소시킨다.
\end{choices}
\end{question}
\answer{3}
